%表紙
\newpage
\chapter{学期について}
\begin{enumerate}[label=第\arabic*条　]
  \item この章は、音楽ゲーム学園学園憲章（以下「憲章」という。）の趣旨に基づき、学園の活動を円滑に行うため、活動期間の単位として「学期」の概念を定め、これに関し必要な事項を定めるものとする。
  \item この章において、次の各号に掲げる用語の意義は、当該各号に定めるところによる。
      \begin{enumerate}[label=(\arabic*)　]
        \item 学期　学園における活動を区分する単位であって、次条に定める正規活動期間及び補講期間を合わせた、半年を一区切りとする期間をいう。
        \item 正規活動期間　学園の講義その他の活動を通常の形態で実施する期間をいう。
        \item 補講期間　正規活動期間の終了後、当該学期における講義の補完的な実施のために設ける期間をいう。
        \item 募集期間　次の学期に向けて、新たに学生となろうとする者の募集を行う期間をいう。
      \end{enumerate}
  \item 学園における学期は、次の各号に掲げる期間の構成により、半年ごとに繰り返し設定する。
      \begin{enumerate}[label=(\arabic*)　]
        \item ４月１日を起点とする学期　正規活動期間を４月１日から７月３１日までとし、補講期間を８月１日から８月１５日までとする。
        \item １０月１日を起点とする学期　正規活動期間を１０月１日から翌年１月３１日までとし、補講期間を翌年２月１日から翌年２月１５日までとする。
      \end{enumerate}
  \item 次の各号に掲げる期間を、それぞれ当該各号に定める学期の募集期間とする。
      \begin{enumerate}[label=(\arabic*)　]
        \item 前条第２号の学期の募集期間　８月１５日から１０月１５日まで
        \item 前条第１号の学期の募集期間　２月１５日から４月１５日まで
      \end{enumerate}
    \begin{enumerate}[label=\arabic*　,start=2,leftmargin=*]
      \item 募集期間は、当該学期の正規活動期間の開始後も継続することができる。
    \end{enumerate}
  \item 学園における学期は、開校（本開講）時における最初の学期を「＃１期」とし、以後、第３条各号の別を問わず、開講の順に１を加えて付番する。
    \begin{enumerate}[label=\arabic*　,start=2,leftmargin=*]
      \item 前項の規定にかかわらず、最初の学期の開始時期その他必要な事項は、学園長が別に定める。
    \end{enumerate}
  \item 学園長は、運営上やむを得ない事由があると認めるときは、第３条及び第４条に定める期間を変更することができる。
    \begin{enumerate}[label=\arabic*　,start=2,leftmargin=*]
      \item 前項により期間を変更したときは、学園長は、速やかにこれを学生に公示しなければならない。
    \end{enumerate}
  \item この章の実施に関し必要な事項は、学園長が別に定める。
\end{enumerate}